%! Interpreting p-Values — Unit 6, Lesson 6.5 — Warm-Up
\documentclass[10pt]{article}
\usepackage{apstats-article}
\usepackage{apstats-boxes}
\newcommand{\TallMath}[1]{$\displaystyle #1\rule[-1.4em]{0pt}{3.2em}$}

\begin{document}

\pageheader{Unit 6, Lesson 6.5}{Warm-Up}
\namedateperiod

\vspace{0.15in}

A quality control engineer claims that 10\% of widgets from a production line are defective.
A random sample of 150 widgets finds 22 defective.

\begin{enumerate}
  \item State appropriate hypotheses (the engineer suspects the defect rate has changed).
        $H_0$: \blank{4cm} \quad $H_a$: \blank{4cm}

  \item Verify conditions.
        \begin{itemize}
          \item Random: \blank{5cm}
          \item Large Counts: \blank{7cm}
          \item 10\%: \blank{5cm}
        \end{itemize}

  \item Compute the test statistic.  ($\hat p = $ \blank{2cm})

        $z = \dfrac{\hat p - p_0}{\sqrt{p_0(1-p_0)/n}} = \dfrac{\blank{1.5cm} - \blank{1.5cm}}{\sqrt{\blank{2.5cm}/\blank{1.5cm}}} \approx$ \blank{2cm}

  \item Looking at the $z$ you computed: should the $p$-value be based on the left tail,
        the right tail, or both tails? Explain.
        \writelines{2}
\end{enumerate}

\end{document}
